% Generated draft: requires mathematical and source review.
\documentclass[11pt]{article}
\usepackage[margin=25mm]{geometry}
\usepackage{fontspec}
\setmainfont{Latin Modern Roman}
\usepackage{amsmath,amssymb,mathtools,unicode-math}
\setmathfont{Latin Modern Math}
\usepackage{enumitem,needspace,xurl,hyperref}
\setlength{\parindent}{0pt}
\setlength{\parskip}{5pt}
\setlength{\emergencystretch}{3em}
\newcommand{\N}{\mathbb{N}}
\newcommand{\Z}{\mathbb{Z}}
\newcommand{\Q}{\mathbb{Q}}
\newcommand{\R}{\mathbb{R}}
\newcommand{\C}{\mathbb{C}}
\newcommand{\F}{\mathbb{F}}
\newcommand{\A}{\mathbb{A}}
\newcommand{\PP}{\mathbb{P}}
\newcommand{\E}{\mathbb{E}}
\newcommand{\eps}{\varepsilon}
\newcommand{\dd}{\,\mathrm{d}}
\begin{document}
\section*{UnsolvedMath ID: 2524}
\textbf{Source problem code:} catalogue entry
\subsection{Definitions and mathematical statement}
For an infinite set $\Omega$, let $S_\Omega$ be the group of all bijections $\Omega\to\Omega$ under composition. Let $B\leq S_\Omega$ be any subgroup, with no closure or measurability assumption. Take two abstract copies of $S_\Omega$ and identify their copies of each $b\in B$. Their amalgamated free product
\[
P=S_\Omega*_B S_\Omega
\]
is the quotient of their ordinary free product by the normal subgroup generated by the relations identifying these two copies of $b$, for all $b\in B$. Its canonical factor maps $j_1,j_2:S_\Omega\to P$ are injective and have intersection exactly the common copy of $B$.

The question asks whether
\[
\forall\text{ infinite sets }\Omega\ \forall B\leq S_\Omega\ \exists\text{ injective group homomorphism }
\varphi:S_\Omega*_B S_\Omega\longrightarrow S_\Omega.
\]
This is abstract group embeddability. The map is not required to restrict to the given inclusion of $B$ into $S_\Omega$, and no topological condition is imposed on it in the question.

The supplied solution answers no already when $\Omega$ is countably infinite. In that case $S_\Omega$ carries the pointwise-convergence topology, whose identity neighborhoods fix finite subsets pointwise; it is a Polish group. Choose a subgroup $B$ that is not Borel in this topology. The source attributes to Corson the following counterexample argument. If an embedding $\varphi$ existed, its two factor restrictions $\varphi_i=\varphi\circ j_i$ would be continuous by the cited automatic-continuity result. As continuous injective images of a Polish space, their images would be Borel by the Lusin--Souslin theorem. The amalgam normal-form intersection property gives
\[
\varphi(j_1(B))=\operatorname{im}(\varphi_1)\cap\operatorname{im}(\varphi_2),
\]
so taking the inverse image under $\varphi_1$ would make $B$ Borel, a contradiction. This is the supplied proof summary, not a newly checked result or a construction for every subgroup.

The background says the question weakens the group case of Bergman's 2007 Question 4.4 and that an earlier 2009 result obstructed embeddings required to respect $B$. The current counterexample is stronger in that it allows any abstract embedding and still rules it out for the chosen non-Borel subgroup. A countably infinite instance suffices to refute the universal question over all infinite sets; no conclusion for every uncountable set or every Borel subgroup is inferred.

The registry labels this problem solved. Its literature triage dated 2026-08-17 agrees, attributes the counterexample to Corson's letter of 26 January 2026 and linked solution, and reports no remaining issue for the yes/no question. These mathematical and status statements are attributed to the supplied sources; no external proof was visited or independently verified here.
\subsection{Short English statement}
Must two copies of an infinite symmetric group, amalgamated over an arbitrary subgroup, embed back into that symmetric group? The supplied solution says no: a non-Borel subgroup in the countable case obstructs even an arbitrary abstract embedding via automatic continuity.
\subsection{Sources}
\begin{itemize}
\item Kourovka Notebook, New Problems, 21st issue (2026), original Problem 21.15; supplied discussion attributed to G. M. Bergman.
\item G. M. Bergman, Indagationes Mathematicae 18 (2007), 349--403, Question 4.4; title not supplied in the record.
\item Related $B$-respecting embedding result, Algebra and Number Theory 3 (2009), 847--879, Section 10, as supplied in the Bergman discussion without a title.
\item Samuel M. Corson, ``Solution of Kourovka Notebook Problem 21.15'' (2026), with letter dated 26 January 2026 in the source: \url{https://kourovkanotebookorg.wordpress.com/wp-content/uploads/2026/03/solution-of-21.15.pdf}. This citation is retained without opening or generating a PDF.
\item Alexander S. Kechris and Christian Rosendal, ``Turbulence, amalgamation, and generic automorphisms of homogeneous structures,'' Proceedings of the London Mathematical Society 94, no. 2 (2007), 302--350: \url{https://arxiv.org/abs/math/0409567}.
\item Local registry data/unsolvedmath/problems.json, permanent catalogue entry, source code catalogue entry, ``Kourovka Notebook Problem 21.15''; statement, complete supplied solution discussion, and literature triage dated 2026-08-17. Dataset attribution: Ulam AI and the UnsolvedMath Contributors.
\end{itemize}
\end{document}
